% Vietnamese translation for OI Public Library
% Source-PDF: 动态规划 DYNAMIC PROGRAMMING/动态规划中的提纯 - 陈启峰.pdf
\documentclass[11pt]{article}
\usepackage{oipl-vn}

\title{``Tinh lọc'' trong quy hoạch động}
\author{Trần Khải Phong}
\date{}

\begin{document}
\maketitle

\origtitle{Purification in Dynamic Programming}
\sourcepath{Dynamic Programming / DP purification - Chen Qifeng.pdf}

\section*{Khái quát}

Trong quy hoạch động có rất nhiều phương pháp tối ưu. Nhìn tổng thể, chúng có
thể chia thành hai loại: \term{giảm tính toán lặp} và \term{giảm dư thừa}. Hai
cách ``tinh lọc'' này đều thể hiện bản chất của quy hoạch động. Giảm tính toán
lặp thường bắt đầu từ chiến lược: trước hết tìm những giá trị bị tính đi tính
lại, sau đó thêm trạng thái mới để lưu giá trị ấy. Còn giảm dư thừa thường có
thể tối ưu từ trạng thái và chiến lược, bằng cách xóa các trạng thái và chiến
lược vô dụng.

Thông qua phần phân tích bài toán dưới đây, hy vọng có thể gợi mở thêm một số
ý tưởng cho người đọc.

\section*{Từ khóa}

Giảm tính toán lặp; giảm dư thừa.

\section{Mô tả bài toán}

Ta cần tìm dãy con chung tăng dài nhất của hai dãy, tức bài toán LCIS.

Cho hai dãy $A$ và $B$ có độ dài lần lượt là $n$ và $m$. Hãy tìm một dãy $C$
dài nhất thỏa mãn:
\[
  C_1<C_2<C_3<\cdots<C_l;
\]
và $C$ vừa là dãy con của $A$, vừa là dãy con của $B$.

\section{Phân tích bài toán}

Khi nhìn thấy ba từ khóa của bài này: dài nhất, chung, tăng, ta rất dễ có cảm
giác quen thuộc. Nếu bỏ yêu cầu ``chung'' hoặc bỏ yêu cầu ``tăng'', ta sẽ nhận
được hai bài toán quen thuộc: dãy con chung dài nhất (LCS) và dãy con tăng dài
nhất (LIS).

Nhiệm vụ đầu tiên của quy hoạch động là xác định trạng thái. Nhìn lại cách
biểu diễn trạng thái của LCS và LIS, rồi kết hợp một cách sáng tạo đặc trưng
của hai cách biểu diễn ấy, ta thu được một dạng trạng thái cho bài này.

Gọi $x_{i,j}$ là độ dài LCIS của $A_1,\ldots,A_i$ và $B_1,\ldots,B_j$, trong
đó yêu cầu $A_i=B_j$ và bắt buộc chọn $A_i$ cùng $B_j$ làm một cặp khớp. Khi
đó ta có phương trình chuyển trạng thái A:
\[
  x_{i,j}=
  \max\left\{
    1,\;
    \max_{\substack{0<k<i,\;0<l<j\\ B_l<B_j\\ x_{k,l}\text{ có nghĩa}}}
      \{x_{k,l}+1\}
  \right\}.
\]
Nếu $A_i\ne B_j$, trạng thái $x_{i,j}$ không có nghĩa và không cần được xét.

Độ phức tạp thời gian của quy hoạch động này là $O(n^4)$, rõ ràng chưa đạt yêu
cầu; ta cần tiến hành ``tinh lọc''.

\section{Cách tinh lọc 1: giảm tính toán lặp}

Phân tích kỹ chiến lược của các trạng thái khác nhau, ta dễ thấy nhiều giá trị
có cùng ý nghĩa bị tính nhiều lần. Chẳng hạn, với hai trạng thái bất kỳ
$x_{i,j_1}$ và $x_{i,j_2}$, nếu tồn tại một chỉ số $l$ sao cho
$B_l<B_{j_1}$, $B_l<B_{j_2}$, đồng thời $l<j_1$ và $l<j_2$, thì khi chọn
chiến lược cho hai trạng thái này, tập trạng thái
\[
  \{x_{k,l}\mid k<i\}
\]
đều bị tính một lần.

Tính một tập trạng thái, về bản chất, là tìm trạng thái có giá trị lớn nhất
trong tập đó. Vì vậy ta chỉ cần biết giá trị lớn nhất của tập là đủ. Các tập
có dạng $\{x_{k,l}\mid k<i\}$ xuất hiện nhiều lần, nghĩa là giá trị lớn nhất
của chúng bị tính lặp. Do đó ta thêm một trạng thái:
\[
  f_{i,j}=\max\{x_{k,j}\mid k\le i\}.
\]
Từ đó thu được phương trình chuyển trạng thái hai lớp B:
\[
  x_{i,j}=
  \max\left\{
    1,\;
    \max_{\substack{0<k<j\\ B_k<B_j}}\{f_{i-1,k}+1\}
  \right\},
\]
\[
  f_{i,j}=\max\{0,\;f_{i-1,j},\;x_{i,j}\}.
\]

Như vậy độ phức tạp thời gian tổng thể là $O(n^3)$. Nếu dùng cây tìm kiếm nhị
phân, độ phức tạp còn có thể giảm xuống $O(n^2\log n)$. Rõ ràng bước ``tinh
lọc'' này đã loại bỏ được rất nhiều tạp chất, nhưng đừng quên ta vẫn còn một
lưỡi dao sắc khác.

\section{Cách tinh lọc 2: giảm dư thừa}

Dư thừa thường ẩn rất kín, khó phát hiện, nên cần dành thêm thời gian và công
sức để đào sâu.

Trong phương trình chuyển trạng thái B, nếu tập trung vào phạm vi chọn của tất
cả chiến lược tối ưu và thử nghiệm nhiều lần, ta sẽ thấy một hiện tượng không
khó nhận ra. Với một trạng thái bất kỳ $x_{i,j}$, nếu chiến lược tối ưu của nó
nhỏ hơn một giá trị cố định
\[
  \operatorname{last}_j=\max\{0,\;w\mid w<j\text{ và }B_w=B_j\},
\]
thì
\[
  x_{i,j}=x_{i,\operatorname{last}_j}.
\]
Từ đó nảy sinh một phỏng đoán: phải chăng với mọi trạng thái $x_{i,j}$, nếu
chiến lược tối ưu của nó nhỏ hơn $\operatorname{last}_j$, thì đều có
$x_{i,j}=x_{i,\operatorname{last}_j}$? Ta chứng minh phỏng đoán này dưới đây.

\paragraph{Chứng minh.}
Với một trạng thái bất kỳ $x_{i,j}$, giá trị $x_{i,\operatorname{last}_j}$ là
một lời giải khả thi. Thật vậy, từ phương án tối ưu của
$x_{i,\operatorname{last}_j}$, ta chỉ cần thay cặp khớp
$(i,\operatorname{last}_j)$ bằng $(i,j)$. Vì
$B_{\operatorname{last}_j}=B_j$, phép thay này giữ nguyên giá trị được khớp và
không phá vỡ tính tăng của dãy. Do đó
\[
  x_{i,j}\ge x_{i,\operatorname{last}_j}.
\]

Tiếp theo, với một $f_{i-1,k}$ bất kỳ trong đó $k<\operatorname{last}_j$, giả
sử phương án tương ứng của $f_{i-1,k}$ có $g$ cặp khớp. Mỗi phương án như vậy
tương ứng một-một với một phương án có $g+1$ cặp khớp của
$x_{i,\operatorname{last}_j}$: chỉ cần thêm cặp khớp
$(i,\operatorname{last}_j)$ vào cuối. Vì thế
\[
  x_{i,\operatorname{last}_j}\ge f_{i-1,k}+1
  \qquad (k<\operatorname{last}_j).
\]

Tổng hợp hai kết luận trên, việc xét các giá trị $f_{i-1,k}+1$ với
$k<\operatorname{last}_j$ là tương đương với việc xét
$x_{i,\operatorname{last}_j}$. Nói cách khác, với mọi trạng thái $x_{i,j}$,
nếu chiến lược tối ưu của nó nhỏ hơn $\operatorname{last}_j$, thì đều có
\[
  x_{i,j}=x_{i,\operatorname{last}_j}.
\]

Sau chứng minh trên, ta cải tiến phương trình chuyển trạng thái B để nhận được
phương trình chuyển trạng thái tốt hơn, ký hiệu C. Quy ước
$x_{i,0}=0$ khi $\operatorname{last}_j=0$:
\[
  x_{i,j}=
  \max\left\{
    1,\;
    \max_{\substack{\operatorname{last}_j<k<j\\ B_k<B_j}}
      \{f_{i-1,k}+1\},\;
    x_{i,\operatorname{last}_j}
  \right\},
\]
\[
  f_{i,j}=\max\{0,\;f_{i-1,j},\;x_{i,j}\},
\]
\[
  \operatorname{last}_j=\max\{0,\;w\mid w<j\text{ và }B_w=B_j\}.
\]

Độ phức tạp thời gian tổng thể của quy hoạch động này chỉ còn $O(n^2)$, thấp
hơn rất nhiều so với trước.

\section*{Tổng kết}

Việc ``tinh lọc'' quy hoạch động không có một hình thức cố định; nó thường cần
một kiểu linh cảm. Linh cảm ấy không phải lâu đài trên không, mà đến từ suy
nghĩ và tích lũy hằng ngày. Chỉ cần thường xuyên tổng kết, thường xuyên suy
nghĩ, thường xuyên đổi mới và kiên trì, ta nhất định có thể học tốt quy hoạch
động.

\section*{Lời cảm ơn}

Xin cảm ơn thầy Tống Tân Ba và Quách Hoa Dương đã đưa ra nhiều ý kiến quý báu
cho bài viết này.

\section*{Tài liệu tham khảo}

\begin{enumerate}
  \item ``Tận dụng đầy đủ tính chất của bài toán: phân tích ví dụ về tối ưu
  hóa `cá thể hóa' trong quy hoạch động'', Hạng Vinh Cảnh, luận văn năm 2003.
  \item ``Giảm dư thừa và tối ưu hóa thuật toán'', Hồ Vĩ Đống, luận văn năm
  2004.
\end{enumerate}

\end{document}
